\section{Appendix C: Code Repository}
\label{appendix:code}

\subsection{Repository Structure}

Miniforge is available at \url{https://github.com/Zapdev-labs/miniforge}.

\begin{verbatim}
miniforge/
|-- src/miniforge/           # Main package
|   |-- __init__.py
|   |-- core/               # Core engine
|   |-- models/             # Model management
|   |-- generation/         # Generation utilities
|   |-- multimodal/         # Vision support
|   -- utils/              # Utilities
|-- benchmarks/             # Benchmark suite
|-- tests/                  # Test suite
|-- examples/               # Usage examples
|-- paper/                  # This paper
|-- configs/                # Sample configurations
|-- README.md
|-- pyproject.toml
-- LICENSE
\end{verbatim}

\subsection{Installation}

\subsubsection{Quick Install}

\begin{verbatim}
pip install miniforge
\end{verbatim}

\subsubsection{Development Install}

\begin{verbatim}
git clone https://github.com/Zapdev-labs/miniforge.git
cd miniforge
uv pip install -e ".[all]"
\end{verbatim}

\subsection{Quick Start Example}

\begin{lstlisting}[language=Python, caption=Basic Usage]
import asyncio
from miniforge import Miniforge

async def main():
    # Load model
    model = await Miniforge.from_pretrained(
        "MiniMaxAI/MiniMax-M2.7",
        quantization="Q4_K_M",
    )
    
    # Simple chat
    response = await model.chat(
        "Explain quantum computing",
        system_prompt="You are a helpful assistant.",
    )
    print(response)
    
    # Streaming
    stream = await model.chat(
        "Tell me a story",
        stream=True
    )
    async for token in stream:
        print(token, end="", flush=True)
    
    # Cleanup
    await model.cleanup()

asyncio.run(main())
\end{lstlisting}

\subsection{Running Benchmarks}

\begin{verbatim}
# Run all benchmarks
python -m benchmarks.runner

# Run specific benchmarks
python -m benchmarks.runner --benchmarks performance memory

# Generate HTML report
python -m benchmarks.runner --html-only --report results/benchmark_*.json
\end{verbatim}

\subsection{License}

Miniforge is released under the MIT License:

\begin{verbatim}
MIT License

Copyright (c) 2026 Zapdev Labs

Permission is hereby granted, free of charge, to any person obtaining a copy
of this software and associated documentation files (the "Software"), to deal
in the Software without restriction, including without limitation the rights
to use, copy, modify, merge, publish, distribute, sublicense, and/or sell
copies of the Software, and to permit persons to whom the Software is
furnished to do so, subject to the following conditions:

The above copyright notice and this permission notice shall be included in all
copies or substantial portions of the Software.

THE SOFTWARE IS PROVIDED "AS IS", WITHOUT WARRANTY OF ANY KIND, EXPRESS OR
IMPLIED, INCLUDING BUT NOT LIMITED TO THE WARRANTIES OF MERCHANTABILITY,
FITNESS FOR A PARTICULAR PURPOSE AND NONINFRINGEMENT.
\end{verbatim}

\subsection{Contributing}

Contributions are welcome! Areas of particular interest:

\begin{itemize}
    \item Additional benchmark tasks
    \item Hardware-specific optimizations
    \item Bug fixes and performance improvements
    \item Documentation improvements
    \item New backend integrations
\end{itemize}

\subsection{Acknowledgments}

Miniforge builds upon several open-source projects:

\begin{itemize}
    \item llama.cpp by Georgi Gerganov
    \item Transformers by HuggingFace
    \item llama-cpp-python by Andrei Betlen
    \item MiniMax model by MiniMax AI
\end{itemize}

\subsection{Contact}

For questions or issues:

\begin{itemize}
    \item GitHub Issues: \url{https://github.com/Zapdev-labs/miniforge/issues}
    \item Email: jacksonwheeler@zapdev.link
\end{itemize}
